\documentclass[11pt]{article}
\usepackage[margin=1in]{geometry}
\usepackage[T1]{fontenc}
\usepackage{lmodern}
\usepackage{amsmath,amssymb,amsthm,mathtools,booktabs,array}
\usepackage{xcolor}
\usepackage{hyperref}
\hypersetup{colorlinks=true,linkcolor=blue!50!black,urlcolor=blue!50!black,
  citecolor=blue!50!black,pdfauthor={Joshua Mark},
  pdftitle={Quantitative refinements and effective experiments for Hadamard existence}}
\setlength{\emergencystretch}{2em}
\newtheorem{theorem}{Theorem}[section]
\newtheorem{proposition}[theorem]{Proposition}
\newtheorem{lemma}[theorem]{Lemma}
\newtheorem{corollary}[theorem]{Corollary}
\newtheorem{remark}[theorem]{Remark}
\newcommand{\NN}{\mathbb N_0}
\newcommand{\HH}{\mathcal H}
\newcommand{\ph}{\Phi}
\newcommand{\ee}{\mathrm e}
\title{Quantitative refinements and effective experiments\\for Hadamard existence}
\author{Joshua Mark\thanks{ORCID: \href{https://orcid.org/0009-0003-5905-2382}{0009-0003-5905-2382}. AI tools assisted with mathematical exploration, formalization, computation and exposition.}}
\date{4 November 2026}
\begin{document}
\maketitle
\begin{abstract}
We refine several consequences of the Hadamard exponent bound
$\tau(p)=O(\log p/\log\log p)$ and examine effective constructions. We derive a uniform flatness
bound with a subpolynomial correction and the condition-number estimate
\[\kappa(n)-1\le n^{-1/4}\exp(O(\log n/\log\log n)).\]
For an independent $r$-generator power-product monoid, a simplex count and
multiset entropy give the second-order upper bound
\[\log\ell_m\le(r-1)m\log m+rm\log\log m+(\beta+o(1))m,\]
with an explicit
constant $\beta$. We also record a quantitative maximal-determinant
consequence and exact finite covering computations. These are deductions
and extensions of the existing result, not a claim of a better order of
growth for $\tau$ or of bibliographic priority.
\end{abstract}

\section{Baseline, domains and status}
For odd positive $p$, let $\tau(p)$ be the least nonnegative integer $t$
for which a Hadamard matrix of order $2^tp$ exists. Write
\[
\ph(x)=\frac{\log x}{\log\log x}\qquad(x>\ee).
\]
All asymptotic statements concern sufficiently large positive integer
inputs. Our existing paper establishes
\begin{equation}\label{eq:baseline}
\tau(p)=O(\ph(p)),\qquad
\exists D>0:\ \forall n\text{ sufficiently large},\quad
\exists d\in\NN:\ d\le g_D(n):=\exp(D\ph(n)),
\end{equation}
where, for every natural $y\le n$, there is $h\in\HH$ with
$y<h\le y+d$. Here $\HH$ is the set of Hadamard orders. The constants
and thresholds are uniform. The proof uses sparse power-product
representations~\cite{KSW} and complementary-sequence transfers
\cite{Craigen,Livinskyi}.

The source conversation's main theorem, covering growth order, gap
theorem and fixed-$\delta$ conditioning corollary are already included
in the current manuscript and Lean project. Repeating them would not be
an improvement. We investigate refinements of the consequences and
effective statements that can be certified by finite computation.
Every assertion below is either proved here, explicitly attributed to a
published theorem, or labelled as a research direction. Mathematical
novelty has not been established by a comprehensive citation review.

\section{Keeping the quantitative scale in flatness and conditioning}
Let $u(n)=\min_{Q\in O(n)}\max_{i,j}|Q_{ij}|$ and let $\kappa(n)$ be the
smallest spectral condition number among sign matrices of order $n$,
with infinity assigned to singular matrices.

\begin{lemma}\label{lem:small}
For fixed $D>0$ and $\varepsilon>0$, $g_D(n)<n^\varepsilon$ eventually.
\end{lemma}
\begin{proof}
The logarithm of $g_D(n)/n^\varepsilon$ is
$\log n\,(D/\log\log n-\varepsilon)$, which tends to $-\infty$.
\end{proof}

\begin{theorem}[Quantitative flatness]\label{thm:flat}
There exist $D>0$ and $N$ such that every $n\ge N$ admits an orthogonal
matrix $Q$ with
\begin{equation}\label{eq:flat}
0<\sqrt n-g_D(n),\qquad
|Q_{ij}|\le\frac1{\sqrt n-g_D(n)}\quad\text{for all }i,j.
\end{equation}
In particular,
\[
u(n)\le\frac1{\sqrt n}\left(1+
O\!\left(n^{-1/2}\exp(D\ph(n))\right)\right).
\]
\end{theorem}
\begin{proof}
Choose $d$ from \eqref{eq:baseline} and apply its gap assertion at $y=n$.
There is a Hadamard order $m>n$ with $k=m-n\le d\le g_D(n)$.
By Lemma~\ref{lem:small}, eventually $g_D(n)<\sqrt n\le\sqrt m$.
The submatrix-repair bound of Alexeev--Jasper--Mixon
\cite[Lemmas 3 and 4]{AJM} constructs an orthogonal matrix of order $n$
whose entries are at most $1/(\sqrt m-k)$ in absolute value.
Since $\sqrt m-k\ge\sqrt n-g_D(n)>0$, \eqref{eq:flat} follows.
The final expansion uses $(1-z)^{-1}=1+O(z)$ with
$z=g_D(n)/\sqrt n\to0$.
\end{proof}

\begin{theorem}[A quantitative endpoint correction]\label{thm:kappa}
There exist $E>0$ and $N$ such that, for all $n\ge N$,
\begin{equation}\label{eq:kappa}
\kappa(n)\le1+n^{-1/4}\exp(E\ph(n)).
\end{equation}
Equivalently the excess has upper bound
$n^{-1/4+O(1/\log\log n)}$. This retains more information than the
family of bounds $1+n^{-1/4+\delta}$ with fixed $\delta>0$.
\end{theorem}
\begin{proof}
Put $g=g_D(n)$ and $v=(\sqrt n-g)^{-1}$. For large $n$, $g\ge1$ and
$g\le\sqrt n/2$, hence $v\le2/\sqrt n$.
Using the flat matrix from Theorem~\ref{thm:flat}, independent sign
rounding has mean $v^{-1}Q$. The quantitative concentration and
singular-value perturbation argument of~\cite[proof of Theorem 1]{AJM}
gives a sample with error at most
\[
e(n,v)=\sqrt{2(n-v^{-2})\log(2n)}+\tfrac23\log(2n).
\]
The variance parameter is nonnegative since $v\ge1/\sqrt n$, and
\[
n-v^{-2}=2\sqrt n\,g-g^2\le2\sqrt n\,g.
\]
Consequently
\[
v e(n,v)=O\!\left(n^{-1/4}\sqrt{g\log n}
                 +n^{-1/2}\log n\right).
\]
It tends to zero by Lemma~\ref{lem:small}. For a rounded matrix $A$,
Weyl's inequalities therefore yield
\[
\operatorname{cond}(A)\le\frac{1+v e(n,v)}{1-v e(n,v)}
\le1+4v e(n,v)
\]
once $v e(n,v)\le1/2$. Since
$\sqrt g=\exp((D/2)\ph(n))$ and
$\log\log n=o(\ph(n))$, logarithmic factors and fixed constants can be
absorbed into $\exp(E\ph(n))$ for any sufficiently large fixed $E$.
This proves \eqref{eq:kappa}.
\end{proof}

\paragraph{A finite bound available in the Lean development.}
The project's moment-concentration proof also supplies a fully explicit
finite substitute for the Bernstein error above:
\[
e_{\rm mom}(n,v)=4(2k_n+1)\sqrt{\max\{1,n-v^{-2}\}},\qquad
k_n=\max\{1,\lceil\log_4(2n)\rceil\}.
\]
Given a gap certificate of size $d\le g<\sqrt n$, the condition
$4v e_{\rm mom}(n,v)\le1$ implies the existence of a sign matrix with
condition number at most $1+4v e_{\rm mom}(n,v)$. This finite statement
and Theorem~\ref{thm:flat} are formalized in the new extension. The
moment bound loses an extra logarithmic factor relative to the displayed
Bernstein bound, which is still absorbed by the same asymptotic scale.

\section{A second-order upper bound for the covering function}
Fix multiplicatively independent integers $b_1,\ldots,b_r>1$, $r\ge2$,
and let $S=\{\prod b_i^{a_i}:a_i\in\NN\}$. Let $\ell_m(S)$ be the first
positive integer not representable using at most $m$ members of $S$,
with repetitions permitted. This section's upper bound does not require
that one base is two; that requirement is used by our covering lower bound.
Put
\[
c_i=\log b_i,\quad C=\sum_i c_i,\quad
A=\frac1{r!\prod_i c_i},\quad
\beta=1+r\log(r-1)-\log\!\left(r!\prod_i\log b_i\right).
\]

\begin{lemma}[Simplex count]\label{lem:simplex}
For $L\ge0$ and $T(e^L)=|S\cap[1,e^L]|$,
\[
\frac{L^r}{r!\prod_i c_i}\le T(e^L)
\le\frac{(L+C)^r}{r!\prod_i c_i}.
\]
Thus $T(N)=A(\log N)^r(1+O(1/\log N))$.
\end{lemma}
\begin{proof}
Independence makes exponent tuples correspond to distinct products.
For each tuple with $\sum c_i a_i\le L$, take the half-open box
$\prod_i[c_i a_i,c_i(a_i+1))$, of volume $\prod_i c_i$.
These boxes have disjoint interiors. Their union contains the positive
simplex $\sum x_i\le L$ (up to boundaries), by taking coordinate floors,
and is contained in $\sum x_i\le L+C$. Compare volumes.
\end{proof}

\begin{theorem}[Second-order entropy bound]\label{thm:entropy}
For every fixed $\eta>0$, all sufficiently large $m$ satisfy
\begin{equation}\label{eq:entropy}
\log\ell_m(S)\le(r-1)m\log m+rm\log\log m+(\beta+\eta)m.
\end{equation}
In particular the normalized second-order remainder has limsup at most
$\beta$.
\end{theorem}
\begin{proof}
Set
\[
L_m=(r-1)m\log m+rm\log\log m+(\beta+\eta/2)m,
\qquad N=\lceil e^{L_m}\rceil.
\]
Then $\log N=L_m+o(1)$ and
\[
\log\log N=\log m+\log\log m+\log(r-1)+o(1).
\]
Every sum at most $N$ uses only members of $S\cap[1,N]$. The number
of multisets of at most $m$ such members is exactly $\binom{T(N)+m}{m}$.
This bounds the number of represented positive integers at most $N$,
even though different multisets can have the same sum.
The estimate $m!\ge(m/e)^m$ gives
\begin{align*}
\log\binom{T(N)+m}{m}
&\le m\bigl(1+\log(T(N)+m)-\log m\bigr)\\
&=(r-1)m\log m+rm\log\log m+(\beta+o(1))m.
\end{align*}
Here $T(N)/m\to\infty$ and Lemma~\ref{lem:simplex} were used in the
last line. This is strictly less than $\log N$ for large $m$, so a
positive integer at most $N$ is missing. Hence $\ell_m(S)\le N$.
The ceiling error $o(1)$ is absorbed by the remaining $\eta m/2$ margin.
\end{proof}

For $S(2,10)$ the result reads
\[
\log\ell_m\le m\log m+2m\log\log m+
\bigl(1-\log(2\log2\log10)+\eta\bigr)m.
\]
For $S(2,10,26)$ it reads
\[
\log\ell_m\le2m\log m+3m\log\log m+
\bigl(1+3\log2-\log(6\log2\log10\log26)+\eta\bigr)m.
\]
The linear constants are approximately $-0.161$ and $-0.361$,
respectively; exact logarithmic expressions, rather than rounded decimals,
are used in the theorem. This refines the existing $(r-1+o(1))m\log m$
upper bound, but does not improve the leading coefficient $r-1$.

\section{A maximal-determinant consequence}
Let $D(n)$ be the largest absolute determinant of a sign matrix of order
$n$, and $R(n)=D(n)/n^{n/2}$.

\begin{theorem}\label{thm:det}
There exists $E>0$ such that eventually
\[
\exp\bigl(-\exp(E\ph(n))\bigr)\le R(n)\le1.
\]
Thus the logarithmic deficit from Hadamard's determinant bound is at most
$\exp(O(\ph(n)))=n^{o(1)}$.
\end{theorem}
\begin{proof}
Take a natural gap certificate $d\le g_D(n)$ as in
\eqref{eq:baseline}; eventually $d<n$. Apply it at $y=n-d$ to find a
Hadamard order in $(n-d,n]$. The largest Hadamard order $h\le n$ therefore
has $a=n-h<d$. If $a=0$, then $R(n)=1$. Otherwise Lemma~\ref{lem:small}
shows that $6a^3/h\to0$, so the hypothesis $h\ge6a^3$ of
Brent--Osborn--Smith~\cite[Theorem 3]{BOS} holds eventually. Their bound
$R(n)\ge0.07(0.352)^a>3^{-(a+3)}$ then applies. Thus we do not need
the finite exceptional-case computations behind their universal Theorem 4.
In either case,
\[
-\log R(n)\le(\log3)(d+3)\le(\log3)(g_D(n)+3)
\le\exp(E\ph(n))
\]
for some fixed $E>D$ and large $n$. Hadamard's determinant inequality
gives the upper bound.
\end{proof}

There is also a short route that avoids the probabilistic determinant
theorem. If $m=n+k$ is a nearby larger Hadamard order, select any $n$
rows. Cauchy--Binet implies that the squared determinants of their
$n\times n$ minors sum to $m^n$. Some minor therefore has determinant
at least $m^{n/2}/\sqrt{\binom{m}{n}}$. In particular
\[
R(n)\ge\frac{(m/n)^{n/2}}{\sqrt{\binom{m}{k}}}
\ge m^{-k/2}.
\]
With $k\le g_D(n)$, the extra factor $\log m$ in the logarithmic
deficit is absorbed by $\exp(E\ph(n))$. This gives the same stated
scale. Neither this elementary minor estimate nor the published
determinant theorem is claimed to be a new ingredient.

\section{Finite covering certificates and their limits}
The exact experiments use bitset subset dynamic programming for all
allowed values at most $2^{22}=4194304$. Descending updates enforce
distinct summands. For a monoid containing two, the resulting at-most-$m$
coverage agrees with the version permitting repetitions: merge two equal
summands into their double until all are distinct, never increasing
the term count. The bitset results were checked against an independent
unbounded-coin dynamic program on smaller intervals.

\begin{center}
\begin{tabular}{r r r r}
\toprule
$m$ & $\ell_m(S(2,3))$ & $\ell_m(S(2,10))$ & $\ell_m(S(2,10,26))$\\
\midrule
1 & 5 & 3 & 3\\
2 & 23 & 7 & 7\\
3 & 431 & 39 & 39\\
4 & 18431 & 199 & 199\\
5 & 3448733 & 1479 & 11999\\
6 & $>4194304$ & 7879 & 637399\\
7 & $>4194304$ & 375495 & $>4194304$\\
\bottomrule
\end{tabular}
\end{center}
An inequality in this table is only a cutoff bound, not a computed first
missing value. The $(2,3)$ column is a numerical benchmark from the
double-base literature; three is not asserted to be a binary Golay length.

\begin{proposition}[Finite radix certificate]\label{prop:radix}
Suppose every integer $0\le a<2^k$ is a sum of at most $w$ binary Golay
lengths, and doubling a permitted length remains permitted. For odd
$p>1$, put $q=(p-1)/2$ and
$d=\lceil\log_{2^k}(q+1)\rceil$. Then a Hadamard matrix exists at some
exponent
\[
t\le2wd+2.
\]
In particular the resulting explicit logarithmic coefficient is $2w/k$.
\end{proposition}
\begin{proof}
Write $q$ in radix $2^k$ with at most $d$ digits. Expand each digit using
its certified Golay-length sum, and multiply its summands by the relevant
power of $2^k$. These are again Golay lengths. There are at most $wd$
summands in total. Craigen's transfer gives $t=2r+2\le2wd+2$.
Repeated lengths are allowed in that transfer, so cross-digit collisions
do not invalidate the construction.
\end{proof}

For example, $\ell_6(S(2,10,26))=637399>2^{19}$ gives $w=6$, $k=19$
and coefficient $12/19\approx0.632$. This is a rigorous explicit
certificate, but it is worse than Du--Jiang's coefficient $3/20$~\cite{DJ}.
Our finite experiment therefore does not establish a competitive uniform
numerical improvement. Larger finite searches alone need not overcome
the asymptotic preprocessing issue.

The separate entropy certificates use integer arithmetic only. For
$N=2^k$, count $T(N)$ exactly by summing
\[
T(N)=\sum_{b,c:\,10^b26^c\le N}
\left(1+\left\lfloor\log_2\left\lfloor
\frac{N}{10^b26^c}\right\rfloor\right\rfloor\right).
\]
If $\binom{T(N)+m}{m}<N$, then $\ell_m\le N$ without enumerating all
integers through $N$. The supplied script records the actual integers
and the strict inequality. These are upper bounds on the first missing
integer, not lower coverage certificates.

\section{What remains a genuine improvement target}
Within a fixed finite independent-base representation system, sparse
weights are bounded below by a positive multiple of $\ph(q)$ for
infinitely many $q$~\cite{KSW}. Thus simply optimizing a representation
inside $S(2,10)$ or $S(2,10,26)$ cannot establish a uniformly smaller
order than $O(\ph(p))$ for the exponent produced by $t=2r+2$.
This is a limitation of that mechanism, not a lower bound on $\tau(p)$.

Three worthwhile open tasks are: obtain usable constants and thresholds
for the logarithm-ratio separation needed by the Golay-specific greedy
proof; exploit complementary sequence families beyond a fixed finite
power-product alphabet; or reduce the matrix-transfer overhead through
additional structure. Du--Jiang also discuss a possible general transfer
from four-phase complementary sets whose resolution could substantially
change their exponent rate~\cite[conclusion]{DJ}. No such new transfer
is proved here.

The source conversation quoted a coefficient $1/5$ from the thesis.
Du--Jiang explicitly flag a stronger thesis claim as depending on an
incorrect earlier theorem. Their safe comparison uses the coefficient
$3/13$ from a different thesis bound~\cite[introduction]{DJ}. None of
our deductions uses the disputed result. Likewise, qualitative
subpolynomial gaps were already attributed to Smith by~\cite{BOS};
the refinements here are not a priority claim for that qualitative fact.

\section{Files and formal verification scope}
All new artifacts are under \path{improvement/}. No proof file, source
manifest or manuscript in the original repository has been changed.
The Lean extension imports the existing project and mathlib; no new
third-party proof code is imported. It proves the quantitative flatness
theorem and a finite error-preserving conditioning theorem. The proofs
use the original unconditional quantitative gap endpoint.
The asymptotic conditioning deduction, the second-order entropy theorem,
and the determinant corollary have complete mathematical proofs above;
this note does not claim Lean formalizations of those three complete
asymptotic statements.

The exact scripts and JSON reports record the finite checks. The
build instructions distinguish kernel-checked extensions, exact finite
computations and paper-only deductions. All project-authored code is
Apache-2.0; mathematical sources are cited rather than copied as code.

\begin{thebibliography}{9}
\raggedright
\bibitem{KSW} D. Krenn, V. Suppakitpaisarn, and S. Wagner,
\emph{On the minimal Hamming weight of a multi-base representation},
Journal of Number Theory 208 (2020), 168--179.
\url{https://arxiv.org/abs/1808.06330}.
\bibitem{Craigen} R. Craigen, \emph{Signed groups, sequences, and the
asymptotic existence of Hadamard matrices}, J. Combin. Theory Ser. A
71 (1995), 241--254. \url{https://doi.org/10.1016/0097-3165(95)90002-0}.
\bibitem{Livinskyi} I. Livinskyi, \emph{Asymptotic Existence of Hadamard
Matrices}, M.Sc. thesis, University of Manitoba, 2012, Theorem 3.10.
\bibitem{DJ} C. Du and Y. Jiang, \emph{Golay Complementary Sequences of
Arbitrary Length and Asymptotic Existence of Hadamard Matrices}, 2024.
\url{https://arxiv.org/abs/2401.15381}.
\bibitem{AJM} B. Alexeev, J. Jasper, and D. G. Mixon,
\emph{Asymptotically optimal approximate Hadamard matrices},
arXiv:2511.14653v2, 2026. \url{https://arxiv.org/abs/2511.14653v2}.
\bibitem{BOS} R. P. Brent, J.-A. H. Osborn, and W. D. Smith,
\emph{Lower bounds on maximal determinants of $\{\pm1\}$ matrices via
the probabilistic method}, arXiv:1211.3248v3, 2013.
\url{https://arxiv.org/abs/1211.3248v3}.
\end{thebibliography}
\end{document}
